\section{Space Complexity --- Codec --- Scan}
%%% generated by the blueprint pipeline from TCSlib.Complexity.SpaceComplexity.Codec.Scan
\subsection*{Space Complexity --- Codec --- Scan: Introduction}

Declarations of the Lean module \texttt{TCSlib.Complexity.SpaceComplexity.Codec.Scan}.

\subsection*{Space Complexity --- Codec --- Scan: Results}

\begin{definition}[Addressed state]
\lean{Complexity.Codec.AdSt}
\label{Complexity.Codec.AdSt}
\leanok
An addressed state consists of a state value $s\in\sigma$, a pointer
$i\in\{0,\ldots,w\}$, and an array assigning an $\alpha$-value to each of the $w+1$
positions.
\end{definition}

\begin{definition}[Addressed cyclic update step]
\lean{Complexity.Codec.adStep}
\label{Complexity.Codec.adStep}
\leanok
\uses{Complexity.Codec.AdSt}
Let $\delta$ map a state and a list of $\alpha$-values to a state, and let an addressed
state have state component $s$, pointer $i$, and array $A$ indexed by $0,\ldots,w$. On
input $a\in\alpha$, first replace $A(i)$ by $a$. If $i=w$, apply $\delta$ to $s$ and the
updated array read in index order as a list, then set the pointer to $0$; otherwise,
keep $s$ unchanged and advance the pointer to $(i+1)\bmod(w+1)$. In either case, the
resulting array is the updated array.
\end{definition}

\begin{lemma}[Processing a partial buffer segment]
\lean{Complexity.Codec.foldl_adStep_partial}
\label{Complexity.Codec.foldl_adStep_partial}
\leanok
\uses{Complexity.Codec.adStep}
\difficulty{6}
Let $\sigma$ and $\alpha$ be types, let $w,j$ be natural numbers, let $\delta : \sigma
\to \mathrm{List}(\alpha) \to \sigma$, and let $l$ be a list over $\alpha$. Suppose
$j+|l|\le w$. Starting from state $s$, position $j$, and any buffer of $w+1$ entries,
folding the update operation over $l$ leaves the state equal to $s$ and advances the
position to $j+|l|$. There is a resulting buffer whose entry at each position $t$ with
$j\le t<j+|l|$ is the entry of $l$ at index $t-j$.
\end{lemma}

\begin{lemma}[Processing one full buffer chunk]
\lean{Complexity.Codec.foldl_adStep_chunk}
\label{Complexity.Codec.foldl_adStep_chunk}
\leanok
\uses{Complexity.Codec.adStep, Complexity.Codec.foldl_adStep_partial}
\difficulty{5}
Let $\sigma$ and $\alpha$ be types, let $w$ be a natural number, and let $\delta$ assign
a state in $\sigma$ to a state and a finite list over $\alpha$. Starting from state $s$,
position zero, and any buffer of $w+1$ entries, folding the update operation over a list
$ch$ of length $w+1$ ends in state $\delta(s,ch)$ at position zero, with some buffer.
\end{lemma}

\begin{theorem}[Processing full chunks of symbols]
\lean{Complexity.Codec.foldl_adStep_flatten}
\label{Complexity.Codec.foldl_adStep_flatten}
\leanok
\uses{Complexity.Codec.adStep, Complexity.Codec.foldl_adStep_chunk}
\difficulty{4}
Let $\sigma$ and $\alpha$ be types, let $w$ be a natural number, let $\delta$ assign an
element of $\sigma$ to each element of $\sigma$ and each finite list of elements of
$\alpha$, and let $s\in\sigma$. Suppose $L$ is a finite list of finite lists of elements
of $\alpha$, each of length $w+1$, and let the initial buffer assign an element of
$\alpha$ to each of the $w+1$ positions. Processing the concatenation of the lists in
$L$ one symbol at a time, updating the state by $\delta$ whenever a full block of $w+1$
symbols has been processed, ends with state equal to the result of applying $\delta$
successively to the lists in $L$, starting from $s$, and with position zero, for some
final buffer.
\end{theorem}

\begin{theorem}[Flattening commutes with aligned zipping]
\lean{Complexity.Codec.zip_flatten}
\label{Complexity.Codec.zip_flatten}
\leanok
\difficulty{4}
Let $L_1$ and $L_2$ be lists of lists, with entries in types $\alpha$ and $\beta$,
respectively. If the two outer lists have equal length and their entries have equal
lengths at every corresponding index, then zipping the flattened lists gives the same
list as zipping corresponding inner lists and flattening the result.
\end{theorem}

\begin{definition}[Three-value state with flag]
\lean{Complexity.Codec.SlSt}
\label{Complexity.Codec.SlSt}
\leanok
For a type $X$, a state consists of three elements of $X$ and one Boolean flag.
\end{definition}

\begin{definition}[Three-value verification step]
\lean{Complexity.Codec.slStep}
\label{Complexity.Codec.slStep}
\leanok
\uses{Complexity.Codec.SlSt}
Let $G\colon X^3\to X$, let the current state have components $(x_1,x_2,x_3,b)$, and let
the input pair be $(u,v)\in X^2$. The next state is
$(x_2,u,v,b\mathbin{\wedge}[x_3=G(x_1,x_2,u)])$, where the bracket is true exactly when
the equality holds.
\end{definition}

\begin{lemma}[Folded three-term recurrence check]
\lean{Complexity.Codec.foldl_slStep_range}
\label{Complexity.Codec.foldl_slStep_range}
\leanok
\uses{Complexity.Codec.slStep}
\difficulty{6}
Let $X$ be a type with decidable equality, let $G:X\times X\times X\to X$, let $e\in X$,
and let $f,g:\mathbb{N}\to X$ and $J\in\mathbb{N}$. Starting from
$(e,e,e,\mathrm{true})$, process the pairs $(f(i),g(i))$ for $0\le i<J$ in order,
updating $(a,b,c,q)$ to $(b,f(i),g(i),q\land[b=G(a,b,f(i))])$. The resulting state is
$(F_J,F_{J+1},H_J,Q_J)$, where $F_k$ is the value at position $k$ in the list
$(e,e,f(0),\ldots,f(J-1))$ (using $e$ when out of range), $H_k$ is the value at position
$k$ in $(e,g(0),\ldots,g(J-1))$ (again using $e$ out of range), and $Q_J$ is true
exactly when $H_i=G(F_i,F_{i+1},F_{i+2})$ for every $i<J$.
\end{lemma}

\begin{theorem}[Sliding recurrence check]
\lean{Complexity.Codec.slide_ok_iff}
\label{Complexity.Codec.slide_ok_iff}
\leanok
\uses{Complexity.Codec.foldl_slStep_range, Complexity.Codec.slStep}
\difficulty{6}
Let $X$ be a set, let $G\colon X^3\to X$, let $e\in X$, and let $u,v$ be finite lists in
$X$ with $|u|=|v|$. The Boolean flag obtained by starting from $(e,e,e,\mathrm{true})$,
processing the pairs in the zip of $u$ and $v$ followed by two copies of $(e,e)$, and at
each step updating the triple and retaining the flag exactly when the previous middle
entry equals $G$ applied to the previous triple, is true if and only if, for every
integer $i<|u|+2$, the entry at position $i$ of the list $e::v$ equals $G$ applied to
the entries at positions $i,i+1,i+2$ of $e::e::u$, where any out-of-range entry is
interpreted as $e$.
\end{theorem}
